%---------------------------------------------------------------------------%
%->> Frontmatter
%---------------------------------------------------------------------------%
%-
%-> 生成封面
%-

\maketitle% 生成中文封面
\MAKETITLE% 生成英文封面
%-
%-> 作者声明
%-
\makedeclaration% 生成声明页
%-
%-> 中文摘要
%-
\intobmk\chapter*{摘\quad 要}% 显示在书签但不显示在目录
\setcounter{page}{1}% 开始页码
\pagenumbering{Roman}% 页码符号

语音是人类最自然的交流方式，语音到语音生成旨在实现从源语音到目标语音的直接转换，涵盖语音翻译和语音对话等典型任务，是构建自然人机交互系统的核心技术。传统的语音到语音生成系统采用级联架构，将语音识别、文本处理和语音合成等独立模块串联执行。级联架构虽然可以复用各模块成熟的预训练模型，但存在错误逐级累积、副语言信息丢失以及串行推理延迟高等固有缺陷。为克服这些局限，研究者们开始探索端到端方法，旨在使用单一模型直接完成从输入语音到输出语音的映射。然而，端到端方法面临两个核心挑战：语音信号的连续性和复杂性使得从源语音到目标语音的直接映射关系难以学习，且语音模态的平行数据远比文本稀缺。

针对上述挑战，本文的核心策略是借助文本生成来引导语音生成——通过先生成目标文本再基于文本表示生成语音的方式进行任务解耦，在降低建模难度的同时追求高质量、低延迟的实时语音交互。本文沿着四个递进的科学问题展开系统性研究，共完成五项具体工作。针对预训练模型的组合与利用问题，提出基于模块化组合的端到端语音到语音翻译模型ComSpeech；针对建模能力与解码效率的矛盾，提出基于有向无环图的语音到语音翻译模型DASpeech；针对多模态序列的同步流式生成问题，提出基于非自回归流式解码的语音交互大模型LLaMA-Omni和基于自回归流式解码的语音交互大模型LLaMA-Omni 2；针对双向语音流的统一序列建模问题，提出基于多信道交织的全双工语音交互模型LLaMA-Omni-Duplex。

\textbf{1. 基于模块化组合的端到端语音到语音翻译模型}

针对预训练模型的组合与利用问题，端到端语音到语音翻译通常需要稀缺的平行语音数据从头训练，无法像两阶段级联系统那样直接复用已有的语音到文本翻译和语音合成预训练模型，其根本原因在于不同模型的词表和表示空间不兼容。本文提出基于模块化组合的端到端语音到语音翻译模型（Composable Speech-to-Speech Translation, ComSpeech）。ComSpeech设计了基于连接时序分类的词表适配器，将语音到文本翻译模型的子词表示转换为语音合成模型的音素表示，打通了两类模型之间的接口壁垒，使任意预训练模型可自由组合为完整的端到端系统。实验结果表明，ComSpeech在有监督场景下超越了所有端到端基线模型，在零样本场景下无需任何平行语音数据仍超越了所有有监督基线。

\textbf{2. 基于有向无环图的语音到语音翻译模型}

针对建模能力与解码效率的矛盾，在端到端可行的基础上，自回归解码面临语音序列过长导致的高延迟问题，而已有非自回归方法虽然解码速度快，但翻译质量与自回归模型存在较大差距。本文提出基于有向无环图的语音到语音翻译模型（Directed Acyclic Transformer for Speech-to-Speech Translation, DASpeech）。DASpeech采用两阶段非自回归解码框架，以有向无环图结构建模目标序列的多种可能路径，有效缓解了非自回归翻译中的多峰性问题。实验结果表明，DASpeech在翻译质量上达到甚至超过自回归模型的水平，同时实现了高达18.53倍的解码加速。

\textbf{3. 基于非自回归流式解码的实时语音交互大模型}

针对多模态序列的同步流式生成问题，进入以大语言模型为核心的仅解码器架构后，语音交互要求在大语言模型生成文本的同时同步产生语音输出，而非自回归解码器延迟低但生成质量不足。本文提出基于非自回归流式解码的实时语音交互大模型LLaMA-Omni。LLaMA-Omni设计了基于连接时序分类的非自回归流式语音解码器，在大语言模型生成文本的同时并行预测离散语音单元，实现文本与语音的同步流式生成。实验结果表明，LLaMA-Omni的响应延迟低至226毫秒，显著优于级联系统和同类基线模型。

\textbf{4. 基于自回归流式解码的高质量语音交互大模型}

进一步地，LLaMA-Omni的非自回归解码器由于各位置预测相互独立，缺乏对语音词元序列依赖关系的建模，导致生成语音的准确性和自然度不足。为在保持低延迟的同时提升语音生成质量，本文提出基于自回归流式解码的高质量语音交互大模型LLaMA-Omni 2。LLaMA-Omni 2设计了基于读写交替策略的自回归流式语音解码器，通过门控融合机制自适应地利用大语言模型的隐状态和文本词元嵌入，在大语言模型生成文本的过程中交替生成对应的语音词元。实验结果表明，LLaMA-Omni 2在保持实时交互的同时大幅提升了语音的内容准确性和自然度。

\textbf{5. 基于多信道交织的全双工语音交互模型}

针对双向语音流的统一序列建模问题，上述语音交互模型均面向基于轮次的对话模式，依赖外部的语音活动检测模块判断用户发言边界，无法应对打断、停顿等复杂交互场景。本文提出基于多信道交织的全双工语音交互模型LLaMA-Omni-Duplex。该模型将用户语音、模型回复文本和模型回复语音三个信道按固定比例组织成交织序列，统一由大语言模型进行自回归建模，通过文本信道中的特殊标记将全双工场景下的时机决策转化为下一词元预测问题，无需额外的分类模块。实验结果表明，LLaMA-Omni-Duplex的轮替成功率达到92\%，打断成功率达到100\%，回复质量显著超越基线模型。

\keywords{语音到语音生成，语音翻译，语音交互，大语言模型，全双工交互}% 中文关键词
%-
%-> 英文摘要
%-
\intobmk\chapter*{Abstract}% 显示在书签但不显示在目录

Speech is the most natural form of human communication. Speech-to-speech generation aims to directly convert source speech into target speech, encompassing tasks such as speech translation and spoken dialogue, and serves as a core technology for building natural human-computer interaction systems. Traditional speech-to-speech generation systems adopt cascaded architectures that chain independent modules such as speech recognition, text processing, and speech synthesis. While cascaded systems can reuse mature pretrained models for each module, they suffer from inherent limitations including error propagation across modules, loss of paralinguistic information, and high latency due to sequential inference. To overcome these limitations, researchers have explored end-to-end approaches that use a single model to directly map input speech to output speech. However, end-to-end approaches face two core challenges: the continuous and complex nature of speech signals makes the direct mapping from source to target speech extremely difficult to learn, and parallel data in the speech modality is far scarcer than in text.

To address these challenges, the core strategy of this thesis is to leverage text generation to guide speech generation---decomposing the task by first generating target text and then generating speech based on the text representation, reducing modeling difficulty while pursuing high-quality, low-latency real-time speech interaction. This thesis conducts systematic research on four progressive scientific problems, resulting in five concrete works. To address the problem of pretrained model composition and utilization, we propose ComSpeech, a modular end-to-end speech-to-speech translation model. To address the contradiction between modeling capacity and decoding efficiency, we propose DASpeech, a directed acyclic graph based speech-to-speech translation model. To address the problem of synchronous streaming generation of multimodal sequences, we propose LLaMA-Omni and LLaMA-Omni 2, two speech interaction models based on non-autoregressive and autoregressive streaming decoding, respectively. To address the problem of unified sequence modeling of bidirectional speech streams, we propose LLaMA-Omni-Duplex, a full-duplex speech interaction model based on multi-channel interleaving.

\textbf{1. Modular End-to-End Speech-to-Speech Translation Model}

To address the problem of pretrained model composition and utilization, end-to-end speech-to-speech translation typically requires scarce parallel speech data for training from scratch and cannot directly reuse pretrained speech-to-text translation (S2TT) and text-to-speech (TTS) models as two-stage cascaded systems do, fundamentally because different models use incompatible vocabularies and representation spaces. We propose ComSpeech (Composable Speech-to-Speech Translation), a modular end-to-end speech-to-speech translation model. ComSpeech designs a CTC-based vocabulary adaptor that converts subword representations from S2TT models into phoneme representations for TTS models, bridging the interface gap between the two types of models and enabling flexible composition of arbitrary pretrained models into a complete end-to-end system. Experimental results show that ComSpeech surpasses all end-to-end baselines in supervised settings and still outperforms all supervised baselines in zero-shot settings without any parallel speech data.

\textbf{2. Directed Acyclic Graph Based Speech-to-Speech Translation Model}

To address the contradiction between modeling capacity and decoding efficiency, building upon the feasibility of end-to-end translation, autoregressive decoding faces high latency due to the excessive length of speech sequences, while existing non-autoregressive methods achieve fast decoding but with significant quality degradation. We propose DASpeech (Directed Acyclic Transformer for Speech-to-Speech Translation). DASpeech adopts a two-stage non-autoregressive decoding framework with directed acyclic graph structure to model multiple possible paths of the target sequence, effectively alleviating the multi-modality problem in non-autoregressive translation. Experimental results show that DASpeech achieves translation quality comparable to or better than autoregressive models with up to 18.53x decoding speedup.

\textbf{3. Real-Time Speech Interaction Model Based on Non-Autoregressive Streaming Decoding}

To address the problem of synchronous streaming generation of multimodal sequences, moving to decoder-only architectures centered on large language models, speech interaction requires synchronous streaming generation of text and speech while the large language model generates text, yet non-autoregressive decoders offer low latency but limited generation quality. We propose LLaMA-Omni, a real-time speech interaction model based on non-autoregressive streaming decoding. LLaMA-Omni designs a CTC-based non-autoregressive streaming speech decoder that predicts discrete speech units in parallel while the large language model generates text, enabling synchronous streaming generation of text and speech. Experimental results show that LLaMA-Omni achieves response latency as low as 226 milliseconds, significantly outperforming cascaded systems and comparable baseline models.

\textbf{4. High-Quality Speech Interaction Model Based on Autoregressive Streaming Decoding}

Furthermore, the non-autoregressive decoder in LLaMA-Omni makes independent predictions at each position without modeling sequential dependencies among speech tokens, leading to insufficient accuracy and naturalness of the generated speech. To improve speech generation quality while maintaining low latency, we propose LLaMA-Omni 2, a high-quality speech interaction model based on autoregressive streaming decoding. LLaMA-Omni 2 designs an autoregressive streaming speech decoder based on a read-write interleaving strategy, adaptively fusing the hidden states and text token embeddings of the large language model through a gating mechanism, and alternately generating corresponding speech tokens during the text generation process. Experimental results show that LLaMA-Omni 2 significantly improves speech accuracy and naturalness while maintaining real-time interaction.

\textbf{5. Full-Duplex Speech Interaction Model Based on Multi-Channel Interleaving}

To address the problem of unified sequence modeling of bidirectional speech streams, the above speech interaction models all operate in turn-based dialogue mode, relying on external voice activity detection modules to determine user utterance boundaries, and cannot handle interruptions, pauses, or other complex interaction scenarios. We propose LLaMA-Omni-Duplex, a full-duplex speech interaction model based on multi-channel interleaving. The model organizes user speech, model response text, and model response speech into an interleaved sequence at a fixed ratio, unified under autoregressive modeling by the large language model, with special tokens in the text channel converting timing decisions in full-duplex scenarios into next-token prediction without additional classification modules. Experimental results show that LLaMA-Omni-Duplex achieves a turn-taking success rate of 92\% and an interruption success rate of 100\%, with response quality significantly surpassing baseline models.

\KEYWORDS{Speech-to-Speech Generation, Speech Translation, Speech Interaction, Large Language Model, Full-Duplex Interaction}% 英文关键词

\pagestyle{enfrontmatterstyle}%
\cleardoublepage\pagestyle{frontmatterstyle}%

%---------------------------------------------------------------------------%
